\providecommand{\Fcal}{\mathcal{F}}
\providecommand{\rH}{\mathrm{H}}

\chapter{Algebraic Iwasawa Comparison}
\label{chap:algebraic-comparison}

% archon:covers MR4372220OnTheAnticyclotomicIwasawaTheoryOfRationalEllipticCurvesAtEisensteinPrimes/Algebraic.lean

This chapter isolates the algebraic comparison lane. The local Selmer formalism is taken as given: here we only compare the residual decomposition
\[
E[p]^{ss}\cong \mathbb F_p(\phi)\oplus \mathbb F_p(\psi)
\]
with the corresponding character Selmer groups, and we record the local Euler-factor bookkeeping that produces the \(\lambda\)-formula for \(\mathfrak X_E\).

Throughout, we write
\[
\mathfrak X_E^S:=\rH^1_{\Fcal_{\rm Gr}^S}(K,M_E)^\vee,
\qquad
\mathfrak X_E:=\rH^1_{\Fcal_{\rm Gr}}(K,M_E)^\vee,
\]
and for each split prime \(w\nmid p\) in \(S\) we use
\[
\mathcal P_w(E)=P_w(\ell^{-1}\gamma_w)\in\Lambda
\]
for the local Euler factor of \(E\). The character-side modules and Euler factors are supplied by the local Selmer chapter.

\section{Selmer groups of \texorpdfstring{\(E\)}{E}}

\begin{proposition}[Residual comparison for \(E[p]\)]
  \label{prop:modp-E}
  \lean{MR4372220.Algebraic.residualComparison}
  \uses{def:iwasawa-algebra, cor:characters}
  % SOURCE: [MR4372220], Prop.~\ref{cormodp}, §Algebraic side, Selmer groups of E (read from references/MR4372220-anticyclotomic-iwasawa-source/Eisenstein.tex)
  % SOURCE QUOTE: "0 \rightarrow \rH^1_{\Fcal_{\rm Gr}^S}(K,M_\phi[p]) \rightarrow \rH^1_{\Fcal_{\rm Gr}^S}(K,M_E[p]) \rightarrow"
  \textit{Source: MR4372220, Selmer groups of E.}
  The short exact sequence \(0 \to M_\phi[p] \to M_E[p] \to M_\psi[p] \to 0\) induces a natural exact sequence
  \[
  0 \to \rH^1_{\Fcal_{\rm Gr}^S}(K,M_\phi[p]) \to \rH^1_{\Fcal_{\rm Gr}^S}(K,M_E[p]) \to
  \rH^1_{\Fcal_{\rm Gr}^S}(K,M_\psi[p]) \to 0.
  \]
  In particular, the imprimitive residual Selmer group of \(E\) is finite, and its \(\mathbb F_p\)-dimension is the sum of the corresponding character dimensions.
\end{proposition}

\begin{proof}
The restriction diagram at \(\bar v\) has exact rows for \(\phi\), \(\psi\), and \(E\). Applying the snake lemma gives the exact sequence on \(S\)-imprimitive residual Selmer groups. The finiteness and dimension statement follows by adding the character-side residual dimensions.
\end{proof}

\begin{proposition}[Torsion and \(\mu=0\) for \(E\)]
  \label{prop:modp-E-tors}
  \lean{MR4372220.Algebraic.muZeroE}
  \uses{prop:modp-E}
  % SOURCE: [MR4372220], Prop.~\ref{propmodp}, §Algebraic side, Selmer groups of E (read from references/MR4372220-anticyclotomic-iwasawa-source/Eisenstein.tex)
  % SOURCE QUOTE: "the modules \mathfrak{X}_E^S and \mathfrak{X}_E are both \Lambda-torsion with \mu=0"
  \textit{Source: MR4372220, Selmer groups of E.}
  The natural map
  \[
  \rH^1_{\Fcal_{\rm Gr}^S}(K,M_E[p]) \simeq \rH^1_{\Fcal_{\rm Gr}^S}(K,M_E)[p]
  \]
  identifies the residual and \(p\)-torsion Selmer groups. Consequently, \(\mathfrak X_E^S\) and \(\mathfrak X_E\) are \(\Lambda\)-torsion, and both have \(\mu\)-invariant zero.
\end{proposition}

\begin{proof}
The assumption on \(\phi|_{G_p}\) forces \(E(K_{\bar v})[p]=0\), so the local \(p\)-torsion comparison at \(\bar v\) is an isomorphism. The finiteness of the residual group then implies \(\mu=0\) for \(\mathfrak X_E^S\), and the quotient map \(\mathfrak X_E^S \twoheadrightarrow \mathfrak X_E\) gives the same conclusion for \(\mathfrak X_E\).
\end{proof}

\begin{corollary}[No finite submodules]
  \label{cor:lambda-E}
  \lean{MR4372220.Algebraic.lambdaE}
  \uses{prop:modp-E-tors, lem:local-E-cofree}
  % SOURCE: [MR4372220], Cor.~\ref{prop426}, §Algebraic side, Selmer groups of E (read from references/MR4372220-anticyclotomic-iwasawa-source/Eisenstein.tex)
  % SOURCE QUOTE: "\mathfrak{X}_E^S has no non-trivial finite \Lambda-submodules"
  \textit{Source: MR4372220, Selmer groups of E.}
  The module \(\mathfrak X_E^S\) has no non-trivial finite \(\Lambda\)-submodules, and
  \[
  \lambda(\mathfrak X_E^S)=\dim_{\mathbb F_p}\rH^1_{\Fcal_{\rm Gr}^S}(K,M_E[p]).
  \]
\end{corollary}

\begin{proof}
The dual local cohomology at \(\bar v\) is \(\Lambda\)-cofree, so the standard Greenberg quotient argument removes finite submodules. Combined with the identification of residual and \(p\)-torsion Selmer groups from the previous proposition, this gives the \(\lambda\)-dimension formula.
\end{proof}

\section{Comparison I: Algebraic Iwasawa invariants}

\begin{theorem}[Algebraic Iwasawa comparison]
  \label{MAINalgside}
  \lean{MR4372220.Algebraic.mainAlgSide}
  \uses{prop:modp-E-tors, cor:lambda-E, prop:lambda-theta}
  % SOURCE: [MR4372220], Thm.~\ref{MAINalgside}, §Algebraic side, Comparison I: Algebraic Iwasawa invariants (read from references/MR4372220-anticyclotomic-iwasawa-source/Eisenstein.tex)
  % SOURCE QUOTE: "\mathfrak{X}_E is \Lambda-torsion with \mu(\mathfrak{X}_E)=0"
  \textit{Source: MR4372220, Comparison I: Algebraic Iwasawa invariants.}
  Assume that \(\phi|_{G_p}\neq \mathds{1},\omega\). Then \(\mathfrak X_E\) is \(\Lambda\)-torsion with \(\mu(\mathfrak X_E)=0\), and
  \[
  \lambda(\mathfrak X_E)=\lambda(\mathfrak X_\phi)+\lambda(\mathfrak X_\psi)
  +\sum_{w\in S}\Bigl(\lambda(\mathcal P_w(\phi))+\lambda(\mathcal P_w(\psi))-\lambda(\mathcal P_w(E))\Bigr).
  \]
\end{theorem}

\begin{proof}
The torsion and \(\mu\)-vanishing part is already supplied by Proposition~\ref{prop:modp-E-tors}. For the \(\lambda\)-invariant, the corollary gives the residual comparison on the \(E\)-side, while the character-side formula supplies the corresponding identity for \(\mathfrak X_\phi\) and \(\mathfrak X_\psi\). The only remaining step is to subtract the local split-prime Euler factors, whose \(E\)-contribution is \(\mathcal P_w(E)\) and whose character contributions are \(\mathcal P_w(\phi)\) and \(\mathcal P_w(\psi)\).
\end{proof}
